\section{热力学第二定律}
\label{sec:热力学第二定律}
在\xref{subsec:热机和热泵}中研究热机和热泵的过程中，我们曾关注到这样的问题
\begin{itemize}
    \item 在热机中我们关注到，热量不可能完全的转化为功，总会在低温热源中损失一些。
    \item 在热泵中我们关注到，热量不可能自发从低温热源向高温热源传输，需要外部做功。
\end{itemize}
包括上述两个例子在内的大量实验事实表明，在自然界中有许多从能量上符合热力学第一定律，但是却并不会自发的进行的过程，这就是说，\empx{自然界的自发过程具有方向性}。而对自发过程方向性的探索，就导致了本节的热力学第二定律的出现，并进一步导致了熵的概念的引入。

\subsection{热力学第二定律}
\label{subsec:热力学第二定律}
热力学第二定律就是从热机和热泵的工作原理，以及如何提高它们的效率和效能的实践中总结出来的，因此，取决于具体对应热泵和热机，热力学第二定律存在两种逻辑上等价的表述。

\begin{OldBoxLaw}[热力学第二定律的克劳修斯表述]
    不可能把热量从低温物体传向高温物体，而不产生其他影响。
\end{OldBoxLaw}

\begin{OldBoxLaw}[热力学第二定律的开尔文表述]
    不可能从单一热源吸收热量，使之完全变为有用功，而不产生其他影响。
\end{OldBoxLaw}
我们将能从单一热源吸收热量并将其完全转化为有用功，且不产生其他影响的循环动作的热机，称为“\uwave{第二类永动机}”。因此开尔文表述也可以表述为：\empx{第二类永动机是不可能实现的}。
\begin{Proof}[克劳修斯表述和开尔文表述的等价性]
\paragraph{若克劳修斯表述不成立，则开尔文表述不成立}
如果克劳修斯表述不成立，那么我们就可以使得热量从低温物体流向高温物体。

这样可以将卡诺热机损耗在低温热源上的热量重新转移至高温热源，恢复低温热源的状态。

这样的净结果就是热机可以仅从高温热源吸热，并将其完全转化为功，违背开尔文表述。

\paragraph{若开尔文表述不成立，则克劳修斯表述不成立}
如果开尔文表述不成立，那么我们就可以使得热量完全转化为功。

这样就可以将卡诺热泵因为外界做功而额外产生的热量重新转化为功，相当于外界未做功。

这样的净结果就是热泵可以在无需外界做功的条件下，将热量逆向传递，违背克劳修斯表述。\vspace{1ex}

综上，克劳修斯表述和开尔文表述是对热力学第二定律的两种等价表述。
\end{Proof}

\subsection{热力学过程的可逆性}
\label{subsec:热力学过程的可逆性}
热力学第二定律两种表述的等价性反映了自然界与热现象有关的宏观过程应当有的一个共同的特征，讨论过程的方向性就是讨论过程的可逆性，由此引出了可逆与不可逆过程的概念。
\begin{OldBoxDefinition}[可逆过程]
    如果热力学系统由某一状态A出发经历某一过程达到另一状态B，那么，如果存在某种过程使得系统从B回到A时，系统的外界也随之恢复原状，或者说，系统在恢复初态的同时消除原过程对外界引起的一切影响，则称其为\uwave{可逆过程}，反之则称为\uwave{不可逆过程}。
\end{OldBoxDefinition}

\empx{可逆过程必须是准静态的}，因为非准静态过程的中间态并不是平衡态，系统没有确定均匀的物态参量，系统在恢复过程中无法重现原有的非均匀过程，典型的例子就是气体对真空的绝热自由膨胀，该过程就是不可逆的，因为气体并不会自发的压缩回另一侧并重新形成真空。

\empx{可逆过程必须是无摩擦的}，因为如果过程中存在摩擦，则系统在正反过程中都要克服摩擦力做功，这些将功将转化为热量，而热量无法完全的转化回功，摩擦带来了能量的耗散。

综合以上的两项的分析，我们可以得出，\empx{只有无摩擦的准静态过程才是可逆的}。

\begin{itemize}
    \item 热力学第二定律的克劳修斯表述指出“热量传递的不可逆性”。
    \item 热力学第二定律的开尔文表述指出“功热转换的不可逆性”。
\end{itemize}

诚然，热机和热泵可以将“功”和“温度差”相互转化，热机可以利用温度差产生功，热泵可以利用功增加温度差，虽然过程中也发生了热量传递和功热转换，但这都是可逆的，然而如果我们直接通过电阻丝将功转化为热，或者直接放任高温物体到低温物体的热量流动。那么在这之后我们想重新从热量中提取功或者重新恢复温度差，且不产生其他影响，就都是不可能的了。

由此可见，热力学第二定律的实质是指出，\empx{一切与热现象有关的实际宏观过程均不可逆}。如果我们从能量的角度看，热力学第二定律描述了能量在转化过程中的自然倾向，能量总是具有转化为热的倾向，热却不能完全的转化为其他形式的能量，这就表明，热或内能虽然也是一种能量，但却更缺乏转化的潜力，或者说，热或内能的能量品质要比其他形式的能量品质更低。

这也或许就是对于未经历物理训练的人而言，能量守恒定律并没有那么显然，因为在每时每刻，高品质的能量都在转化为低品质的热量，汽车在刹车中将动能转化为热，热水壶在烧水时将电能转化为热，而这些热既不能变为动能也不能变为电能，尽管现在我们知道热或者内能也是能量的一种形态，或者说是运动的一种形态，但从朴素的观点看，能量确实是消失了。

这种朴素想法比较确切的表述是，尽管能量守恒，但是能量的品质在不断下降，任何宏观过程的进行都会导致能量的贬值，使可用能量不断减少。这也就是为何能量守恒但会有能源危机。




